\documentclass[../../../include/open-logic-section]{subfiles}

\begin{document}

\olfileid{sth}{spine}{rank}
\olsection{الرتبة}

لما عرّفنا المراحل بمجموعات~$\OLId{latin-looped}{V}_\OLId{greek-variable}{alpha}$، وثبت أن كل مجموعة جزء من
مرحلة ما، أمكن أن نعرّف \emph{رتبة} المجموعة. ورتبة~$\OLId{latin-looped}{A}$ في الحدس
هي أول مرحلة تتكون فيها~$\OLId{latin-looped}{A}$؛ وأما حدها الدقيق فهو الآتي.

\begin{defn}\ollabel{defnsetrank}
لكل مجموعة~$\OLId{latin-looped}{A}$، تكون~$\setrank{\OLId{latin-looped}{A}}$ أصغر عدد ترتيبي~$\OLId{greek-variable}{alpha}$ يحقق
$\OLId{latin-looped}{A} \subseteq \OLId{latin-looped}{V}_\OLId{greek-variable}{alpha}$.
\end{defn}
\begin{prop}\ollabel{ranksexist}
	توجد~$\setrank{\OLId{latin-looped}{A}}$ لكل~$\OLId{latin-looped}{A}$.
\end{prop}
\begin{proof}
	متروك تمرينًا.
\end{proof}
\begin{prob}
	أثبت~\olref[sth][spine][rank]{ranksexist}.
\end{prob}
\noindent
وبفضل الترتيب الجيد للرتب نستخرج نتائج نافعة، أولها:
\begin{prop}\ollabel{valphalowerrank}
لكل عدد ترتيبي~$\OLId{greek-variable}{alpha}$،
$\OLId{latin-looped}{V}_\OLId{greek-variable}{alpha} = \Setabs{\OLId{latin-ordinary}{x}}{\setrank{\OLId{latin-ordinary}{x}} \in \OLId{greek-variable}{alpha}}$.
\end{prop}

\begin{proof}
إذا كان~$\setrank{\OLId{latin-ordinary}{x}} \in \OLId{greek-variable}{alpha}$، كان
$\OLId{latin-ordinary}{x} \subseteq \OLId{latin-looped}{V}_{\setrank{\OLId{latin-ordinary}{x}}} \in \OLId{latin-looped}{V}_\OLId{greek-variable}{alpha}$؛ فيلزم
$\OLId{latin-ordinary}{x} \in \OLId{latin-looped}{V}_\OLId{greek-variable}{alpha}$، لأن~$\OLId{latin-looped}{V}_\OLId{greek-variable}{alpha}$ مقتدرة، باستعمال
\olref[Valphabasic]{Valphabasicprops} مرات. وفي العكس، إذا كان
$\OLId{latin-ordinary}{x} \in \OLId{latin-looped}{V}_\OLId{greek-variable}{alpha}$، كان~$\OLId{latin-ordinary}{x} \subseteq \OLId{latin-looped}{V}_\OLId{greek-variable}{alpha}$، ومن ثم
$\setrank{\OLId{latin-ordinary}{x}} \leq \OLId{greek-variable}{alpha}$. ويبيّن استقراء بسيط متناهٍ في التجاوز أن
$\OLId{latin-ordinary}{x} \notin \OLId{latin-looped}{V}_{\setrank{\OLId{latin-ordinary}{x}}}$؛ ويلزم لذلك~$\setrank{\OLId{latin-ordinary}{x}} \neq \OLId{greek-variable}{alpha}$،
ومن ثم~$\setrank{\OLId{latin-ordinary}{x}} \in \OLId{greek-variable}{alpha}$.
\end{proof}
\begin{prob}
	أكمل الاستقراء البسيط المتناهي التجاوز المذكور في
	\olref[sth][spine][rank]{valphalowerrank}.
\end{prob}

\begin{prop}\ollabel{rankmemberslower}
إذا كان~$\OLId{latin-looped}{B} \in \OLId{latin-looped}{A}$، فإن~$\setrank{\OLId{latin-looped}{B}} \in \setrank{\OLId{latin-looped}{A}}$.
\end{prop}

\begin{proof}
$\OLId{latin-looped}{A} \subseteq \OLId{latin-looped}{V}_{\setrank{\OLId{latin-looped}{A}}} =
\Setabs{\OLId{latin-ordinary}{x}}{\setrank{\OLId{latin-ordinary}{x}} \in \setrank{\OLId{latin-looped}{A}}}$ بموجب
\olref{valphalowerrank}.
\end{proof}
\noindent
ومن هذه الحقيقة نحصل على ضرب من الاستقراء يثبت الأحكام في
\emph{جميع المجموعات}:

\begin{thm}[مخطط الاستقراء بـ$\in$]
لكل صيغة~$\phi$:
\[
	\forall \OLId{latin-looped}{A}((\forall \OLId{latin-ordinary}{x} \in \OLId{latin-looped}{A})\phi(\OLId{latin-ordinary}{x}) \lif \phi(\OLId{latin-looped}{A})) \lif \forall \OLId{latin-looped}{A} \phi(\OLId{latin-looped}{A}).
\]
\end{thm}

\begin{proof}
نبرهن بالمقابل. افترض~$\lnot \forall \OLId{latin-looped}{A} \phi(\OLId{latin-looped}{A})$. يقتضي الاستقراء
المتناهي التجاوز
(\olref[ordinals][basic]{ordinductionschema}) وجود مجموعة لا تحقق
$\phi$ ورتبتها أصغر رتبة ممكنة؛ أي يوجد~$\OLId{latin-looped}{A}$ بحيث~$\lnot \phi(\OLId{latin-looped}{A})$ و
$\forall \OLId{latin-ordinary}{x}(\setrank{\OLId{latin-ordinary}{x}} \in \setrank{\OLId{latin-looped}{A}} \lif \phi(\OLId{latin-ordinary}{x}))$. فإذا كان
$\OLId{latin-ordinary}{x} \in \OLId{latin-looped}{A}$، كان~$\setrank{\OLId{latin-ordinary}{x}} \in \setrank{\OLId{latin-looped}{A}}$ من
\olref{rankmemberslower}، ومنه~$\phi(\OLId{latin-ordinary}{x})$. وهكذا
$(\forall \OLId{latin-ordinary}{x} \in \OLId{latin-looped}{A})\phi(\OLId{latin-ordinary}{x}) \land \lnot \phi(\OLId{latin-looped}{A})$.
\end{proof}\noindent
ومعنى ذلك بغير الرموز أن نقول إن~$\phi$ \emph{وراثية} إذا وفقط إذا
كان تحقق~$\phi$ في كل !!{element} من مجموعة موجبًا لتحقق~$\phi$ في
المجموعة نفسها. فمخطط الاستقراء بـ$\in$ يقول: إذا كانت~$\phi$ وراثية،
فإن كل مجموعة تحقق~$\phi$.

ونختم، إلى حين، بحث الرتب بإثبات دعاوى تقدمت الإشارة إليها.

\begin{prop}\ollabel{ranksupstrict}
$\setrank{\OLId{latin-looped}{A}} = \supstrict_{\OLId{latin-ordinary}{x} \in \OLId{latin-looped}{A}}\setrank{\OLId{latin-ordinary}{x}}$.
\end{prop}

\begin{proof}
ضع~$\OLId{greek-variable}{alpha} = \supstrict_{\OLId{latin-ordinary}{x} \in \OLId{latin-looped}{A}}\setrank{\OLId{latin-ordinary}{x}}$. من
\olref{rankmemberslower} ينتج~$\OLId{greek-variable}{alpha} \leq \setrank{\OLId{latin-looped}{A}}$. ومن جهة
أخرى، إذا كان~$\OLId{latin-ordinary}{x} \in \OLId{latin-looped}{A}$، كان~$\setrank{\OLId{latin-ordinary}{x}} \in \OLId{greek-variable}{alpha}$، ومن ثم
$\OLId{latin-ordinary}{x} \in \OLId{latin-looped}{V}_\OLId{greek-variable}{alpha}$ بموجب~\olref{valphalowerrank}. فيكون
$\OLId{latin-looped}{A} \subseteq \OLId{latin-looped}{V}_\OLId{greek-variable}{alpha}$، أي~$\setrank{\OLId{latin-looped}{A}} \leq \OLId{greek-variable}{alpha}$؛ فتتعين
$\setrank{\OLId{latin-looped}{A}} = \OLId{greek-variable}{alpha}$.
\end{proof}

\begin{cor}\ollabel{ordsetrankalpha}
لكل عدد ترتيبي~$\OLId{greek-variable}{alpha}$، $\setrank{\OLId{greek-variable}{alpha}} = \OLId{greek-variable}{alpha}$.
\end{cor}

\begin{proof}
افترض، في الاستقراء المتناهي التجاوز، أن
$\setrank{\beta} = \beta$ لكل~$\beta \in \OLId{greek-variable}{alpha}$. عندئذ
$\setrank{\OLId{greek-variable}{alpha}} = \supstrict_{\beta \in \OLId{greek-variable}{alpha}}\setrank{\beta}
= \supstrict_{\beta \in \OLId{greek-variable}{alpha}}\beta = \OLId{greek-variable}{alpha}$ بموجب
\olref{ranksupstrict}.
%	First note that $\setrank{\alpha} \neq \beta$ for any $\beta \in
%	\alpha$. For otherwise we would have $\beta = \setrank{\beta} \in
%	\setrank{\alpha} = \beta$ by \olref{rankmemberslower}, i.e.,
%	$\beta \in \beta$, a contradiction. 
%
%	Now note that $\alpha \subseteq V_\alpha$. For $\alpha =
%	\Setabs{\beta}{\beta \in \alpha}$ by
%	\olref[sfr][ordinals][basic]{ordissetofsmallerord}. And if
%	$\beta \in \alpha$ then $\beta \subseteq V_{\setrank{\beta}} =
%	V_\beta \in V_\alpha$, hence $\beta \in V_\alpha$, by
%	\olref[sfr][spine][Valphabasic]{Valphabasicprops} twice. 
\end{proof}

وأخيرًا نثبت سريعًا ما وُعد به في آخر~\olref[foundation]{sec}، وهو أن
$\ZFminus$ تثبت
$\emph{الانتظام} \Rightarrow \emph{التأسيس}$. ومفهوم «الرتبة»
و\olref{rankmemberslower} مباحان في هذا البرهان؛ لأنهما، كما ذكر في
صدر هذا القسم، يُعرضان باستعمال~$\ZFminus + \text{الانتظام}$.

\begin{prop}[العمل في~$\ZFminus + \text{الانتظام}$]\ollabel{zfminusregularityfoundation}
يصدق التأسيس.
\end{prop}

\begin{proof}
ثبّت~$\OLId{latin-looped}{A} \neq \emptyset$، واختر~$\OLId{latin-looped}{B} \in \OLId{latin-looped}{A}$ ذا أصغر رتبة ممكنة. إذا
كان~$\OLId{latin-ordinary}{c} \in \OLId{latin-looped}{B}$، كان~$\setrank{\OLId{latin-ordinary}{c}} \in \setrank{\OLId{latin-looped}{B}}$ بموجب
\olref{rankmemberslower}؛ ولذلك~$\OLId{latin-ordinary}{c} \notin \OLId{latin-looped}{A}$، باختيار~$\OLId{latin-looped}{B}$.
\end{proof}

\end{document}
